\chapter{基于深度射频指纹的复杂环境单目标高精度定位方法}
\label{ch:single_target}

\section{问题描述与总体思路}
\label{sec:3_1}

复杂城市环境中的无线传播具有显著的非理想性。高层建筑密集分布、街道空间狭窄、反射界面丰富，使得低轨卫星相关信号或目标相关辐射信号在传播过程中普遍经历遮挡、反射、绕射和散射。此时，接收端观测到的信号通常不再由单一直达径主导，而是多条传播路径叠加后的结果。传统基于接收信号强度（Received Signal Strength Indicator, RSSI）、到达角（Angle of Arrival, AoA）以及到达时间（Time of Arrival, TOA）和到达时间差（Time Difference of Arrival, TDOA）等低维几何观测量的定位方法，往往隐含了较理想的视距传播假设。当直达路径被阻断或主导分量由反射径构成时，这些观测量与目标几何位置之间的稳定映射关系将被破坏，从而导致定位误差显著增大。

对于 RSSI 而言，接收功率不仅受传播距离影响，还会受到阴影衰落、多径干涉和环境波动的共同作用，因此同一位置下的观测值可能出现明显起伏。AoA 方法依赖于主导来波方向能够代表真实目标方位，而在建筑反射强烈的城市峡谷场景中，接收端获得的主导分量可能对应反射路径而非直达路径。TOA 与 TDOA 类方法虽然在理想同步条件下具有清晰的几何意义，但当首达径被遮挡时，最早到达分量往往已经包含额外传播路径，从而导致距离或距离差观测产生系统性偏移。由此可见，复杂城市环境中的定位困难并非仅来源于噪声增大，而是传播环境本身改变了观测量的物理含义。

尽管多径和非视距传播会削弱传统几何定位方法的有效性，但从另一角度看，多径结构本身也包含了丰富的位置相关信息。不同空间位置对应不同的建筑布局、反射界面组合和传播路径时延结构，这些因素会在接收信号的幅度、相位和时频分布上留下独特的传播特征。如果能够直接从高维接收信号中学习这种由环境塑造的差异，而不是仅依赖少量人工构造的几何观测量，则有望突破传统方法在复杂场景中的性能瓶颈。

基于上述认识，本章提出一种基于深度射频指纹的复杂环境单目标高精度定位方法。该方法将复杂传播环境下的 I/Q 复基带信号视为位置相关特征的主要载体，利用高保真射线追踪技术和真实城市建筑数据构建三维传播场景，并在此基础上生成与位置一一对应的 I/Q 双路样本。随后，采用 ResNet-50 深度残差网络从原始 I/Q 采样中自动提取高维射频指纹特征，并通过分类与回归联合优化完成单目标定位。考虑到实际排查场景中可能出现低信噪比和严重遮挡情形，本章进一步引入多天线空间分集机制，通过多通道输入提升特征稳定性和定位鲁棒性。

本章的技术路线可概括为以下四个步骤：首先，结合 OpenStreetMap（OSM）建筑数据与射线追踪方法构建三维城市传播场景；其次，参考低轨卫星下行链路的典型时频结构构造参考基带信号，并生成包含多径衰落特征的 I/Q 样本；然后，采用 ResNet-50 提取深度射频指纹并建立位置映射；最后，通过 MIMO 空间分集增强低信噪比条件下的特征可分性。图~\ref{fig:ch3_overall_architecture} 给出了本章方法的总体流程。

\begin{figure}[htbp]
    \centering
    % \includegraphics[width=0.85\textwidth]{figures/ch3_workflow.pdf}
    \vspace{4.5cm}
    \caption{基于深度射频指纹的单目标定位方法总体架构}
    \label{fig:ch3_overall_architecture}
\end{figure}

\section{低轨卫星信号特征构造与样本生成}
\label{sec:3_2}

\subsection{Starlink 下行链路帧结构分析}
\label{subsec:3_2_1}

为使信号建模过程与低轨卫星通信场景保持一致，本文选取 Starlink Ku 波段下行链路作为参考建模对象。需要说明的是，本章研究重点在于定位样本构造与特征学习，而非通信协议逆向识别本身。因此，本文并不试图对实际系统协议进行完整复现，而是在参考公开文献已有识别结果的基础上，提取与定位建模直接相关的关键时频参数和同步结构，并在此基础上构造参考基带信号。

公开研究普遍将 Starlink Ku 波段下行链路建模为 OFDM 体制。设单信道带宽为 \(F_s\)，子载波数为 \(N\)，循环前缀长度为 \(N_g\)，则有效符号时长、保护间隔时长和总符号时长分别可表示为
\begin{equation}
T=\frac{N}{F_s},
\label{eq:ofdm_useful_time}
\end{equation}
\begin{equation}
T_g=\frac{N_g}{F_s},
\label{eq:ofdm_guard_time}
\end{equation}
\begin{equation}
T_{\mathrm{sym}}=T+T_g.
\label{eq:ofdm_symbol_time}
\end{equation}
在本文采用的参考参数中，单信道带宽取 \(240\,\mathrm{MHz}\)，子载波数取 \(1024\)，循环前缀长度取 32。由式~\eqref{eq:ofdm_useful_time}--\eqref{eq:ofdm_symbol_time} 可得，有效符号时长约为 \(4.266\,\mu\mathrm{s}\)，保护间隔时长约为 \(0.133\,\mu\mathrm{s}\)，总符号时长约为 \(4.4\,\mu\mathrm{s}\)。此外，参考文献还表明，帧周期约为 \(1/750\,\mathrm{s}\)，每帧包含 302 个非零符号。上述参数共同决定了本文参考信号的时频组织方式。

在频谱组织方面，参考文献指出 Starlink Ku 波段下行链路可视作由多个相邻信道构成的多信道系统，本文采用 8 个信道的近似组织方式，相邻信道中心间隔为 \(250\,\mathrm{MHz}\)，信道间设置 \(10\,\mathrm{MHz}\) 左右的保护带，同时在中心附近保留 4 个子载波宽度的中间保护结构。该设计反映了低轨卫星宽带通信系统在频谱组织上的工程特点，也为本文后续构造多径条件下的基带样本提供了带宽与时频结构基础。

在同步结构方面，本文重点保留主同步字段（PSS）、辅同步字段（SSS）以及少量辅助同步结构。这样做的目的并不是还原实际协议的全部细节，而是利用同步字段良好的相关特性和稳定的帧级时序特征，为后续样本对齐和特征提取提供一致的参考锚点。考虑到低轨卫星高速运动会引入明显的 Doppler 效应，本文在信号建模过程中进一步显式引入频移项，以增强样本的物理一致性。

为避免信号建模参数在正文中分散出现，本文将参考信号构造过程中采用的主要参数统一列于表~\ref{tab:signal_params_ch3}。

\begin{table}[htbp]
    \centering
    \bicaption{低轨卫星下行参考信号建模参数设置}{Modeling parameters for LEO satellite downlink reference signal}
    \label{tab:signal_params_ch3}
    \begin{tabular}{@{}llccl@{}}
        \toprule
        \textbf{参数类别} & \textbf{参数名称} & \textbf{符号} & \textbf{取值/设置} & \textbf{说明} \\ \midrule
        频域参数 & 载波中心频率 & \(f_c\) & Ku 波段相关频率 & 参考信号载频设置 \\
        & 单信道带宽 & \(F_s\) & \(240\,\mathrm{MHz}\) & 单信道有效带宽 \\
        & 子载波数 & \(N\) & 1024 & OFDM 参考参数 \\
        & 循环前缀长度 & \(N_g\) & 32 & OFDM 参考参数 \\
        时域参数 & 有效符号时长 & \(T\) & \(4.266\,\mu\mathrm{s}\) & \(T=N/F_s\) \\
        & 保护间隔时长 & \(T_g\) & \(0.133\,\mu\mathrm{s}\) & \(T_g=N_g/F_s\) \\
        & 总符号时长 & \(T_{\mathrm{sym}}\) & \(4.4\,\mu\mathrm{s}\) & \(T_{\mathrm{sym}}=T+T_g\) \\
        帧参数 & 帧周期 & \(T_f\) & \(1/750\,\mathrm{s}\) & 参考帧级时间结构 \\
        & 每帧非零符号数 & — & 302 & 参考帧级结构参数 \\
        频谱结构 & 信道数量 & — & 8 & 多信道组织方式 \\
        & 相邻信道间隔 & — & \(250\,\mathrm{MHz}\) & 频谱布局参数 \\
        & 保护带宽 & — & \(10\,\mathrm{MHz}\) & 信道间保护带 \\
        & 中心 gutter 宽度 & — & 4 个子载波 & 中心保护结构 \\
        同步结构 & 主同步字段 & PSS & 保留 & 用于同步锚定与特征构造 \\
        & 辅同步字段 & SSS & 保留 & 用于辅助同步与结构对齐 \\
        & 其他辅助字段 & CSS/CM1SS & 按需选取 & 仅保留定位相关结构 \\
        传播参数 & 多普勒频移 & \(f_d\) & 场景相关 & 由星地相对运动关系决定 \\
        信道参数 & 有效传播路径数 & \(L\) & 场景相关 & 由射线追踪结果确定 \\
        噪声参数 & 信噪比范围 & SNR & \([-15,20]\,\mathrm{dB}\) & 用于鲁棒性实验 \\
        样本形式 & 输入样本表示 & \(\mathbf{X}_{\mathrm{in}}\) & I/Q 双通道矩阵 & 用于深度射频指纹学习 \\ \bottomrule
    \end{tabular}
\end{table}

\subsection{I/Q 双路复基带信号重构}
\label{subsec:3_2_2}

在参考基带信号构造完成后，需要进一步将传播场景中的多径参数映射为接收端可处理的 I/Q 复基带样本。传统几何定位方法通常仅利用接收功率、到达角或若干统计特征进行位置估计。这种做法虽然降低了输入维度，但也会丢失大量与传播环境相关的细粒度信息。特别是在复杂城市环境中，多径分量之间的幅度差异、相位旋转以及到达时延结构具有较强的位置辨识性，若在样本构造阶段就对这些信息进行过强压缩，则后续模型难以恢复其与位置之间的非线性映射关系。

基于这一考虑，本文直接采用完整 I/Q 双路复基带采样作为定位模型输入。设发射端参考基带信号为 \(s(t)\)，传播场景中共有 \(L\) 条有效路径。对于第 \(l\) 条路径，射线追踪模块输出其到达方向 \((\theta_l,\varphi_l)\)、传播时延 \(\tau_l\)、Doppler 频移 \(f_{d,l}\) 以及等效复增益 \(\beta_l\)。考虑接收天线方向图增益 \(G(\theta_l,\varphi_l)\) 后，接收端复基带信号可表示为
\begin{equation}
r(t)=\sum_{l=1}^{L}\beta_l\,G(\theta_l,\varphi_l)\,s(t-\tau_l)e^{j2\pi f_{d,l}t}+n(t),
\label{eq:rx_signal_ch3}
\end{equation}
其中，\(n(t)\) 为复加性高斯白噪声。式~\eqref{eq:rx_signal_ch3} 给出了本文样本生成中采用的等效接收信号模型，其中 \(\beta_l\) 已综合反映路径衰减和相位旋转效应。

在数字接收端，以采样率 \(f_s\) 对 \(r(t)\) 进行离散化处理，可得到长度为 \(N_s\) 的复采样序列 \(r[n]\)。为了适配后续卷积网络输入，本文将其拆分为实部和虚部，形成双通道输入矩阵
\begin{equation}
\mathbf{X}_{\mathrm{in}}=
\begin{bmatrix}
\Re\{r[1]\} & \Re\{r[2]\} & \cdots & \Re\{r[N_s]\} \\
\Im\{r[1]\} & \Im\{r[2]\} & \cdots & \Im\{r[N_s]\}
\end{bmatrix}.
\label{eq:iq_matrix_ch3}
\end{equation}
该表示方式能够较完整地保留信号经历复杂传播后的幅度、相位与时序结构，从而为深度射频指纹提取提供高信息密度输入。

需要指出的是，本文并不将 I/Q 样本简单视为通信信号的数值展开，而是将其视为传播环境作用后的综合观测结果。换言之，式~\eqref{eq:iq_matrix_ch3} 中的输入矩阵同时包含通信体制特征与复杂城市场景传播特征，是后续深度指纹学习的直接对象。

\subsection{样本标注与数据预处理}
\label{subsec:3_2_3}

为了训练具有较好泛化能力的深度定位模型，本文基于高保真仿真平台构建大规模射频指纹数据集。首先，利用 OSM 建筑数据建立感兴趣区域的三维城市场景模型，并在该区域内按预定分辨率进行网格化划分。对于每一个网格位置，射线追踪模块输出相应的多径传播参数，再结合式~\eqref{eq:rx_signal_ch3} 和式~\eqref{eq:iq_matrix_ch3} 构造对应的 I/Q 双路样本。

本文最终构建的数据集覆盖目标区域内 27,546 个网格位置。由于本章主要面向二维平面单目标定位问题，回归标签采用平面坐标 \((x_k,y_k)\)。同时，为增强模型对区域级空间结构的判别能力，本文依据网格划分结果为每个样本赋予区域分类标签 \(y_k^{(\mathrm{cls})}\)。因此，后续网络训练包含两类监督信息：一类是区域分类标签，用于学习不同区域之间的判别边界；另一类是位置回归标签，用于学习区域内部更细粒度的连续位置映射。

在数据预处理阶段，为减小不同样本之间的幅值差异并提高训练稳定性，本文对输入张量进行基于通道的标准化处理。设原始输入样本为 \(\mathbf{X}_{\mathrm{in}}\)，则其标准化结果为
\begin{equation}
\tilde{\mathbf{X}}_{\mathrm{in}}^{i,j}=\frac{\mathbf{X}_{\mathrm{in}}^{i,j}-\mu_i}{\sigma_i},
\label{eq:zscore_ch3}
\end{equation}
其中，\(\mu_i\) 与 \(\sigma_i\) 分别表示第 \(i\) 个输入通道在训练集上的均值与标准差。该处理有助于缓解局部异常值和量纲差异对训练过程的干扰，并提升网络收敛稳定性。

在数据划分方面，本文将全部样本按 8:1:1 的比例划分为训练集、验证集和测试集。训练集用于模型参数学习，验证集用于训练过程监控与超参数选择，测试集用于最终性能评估。通过上述划分方式，可以较为客观地评价模型对未知位置样本的泛化能力。

\begin{figure}[htbp]
    \centering
    % \includegraphics[width=0.72\textwidth]{figures/ch3_ray_tracing_urban.pdf}
    \vspace{4.0cm}
    \caption{基于 OSM 数据的三维城市场景建模与射线追踪样本生成示意图}
    \label{fig:ray_tracing_urban}
\end{figure}

\section{基于 ResNet-50 的射频指纹模型设计}
\label{sec:3_3}

\subsection{网络架构设计}
\label{subsec:3_3_1}

复杂城市环境中的 I/Q 采样信号具有明显的高维、非线性和多尺度特征。浅层卷积网络虽然能够提取部分局部波形变化，但难以对跨路径时延、多分量相位耦合以及长时间跨度内的传播模式进行有效建模。另一方面，若仅通过简单增加网络深度来提升表达能力，则容易带来训练退化问题。基于上述考虑，本文选用 ResNet-50 作为射频指纹定位模型的主干网络。

ResNet 的核心思想是在深层网络中引入恒等映射的跳跃连接，使网络学习输入与输出之间的残差关系。设某一残差模块的输入为 \(\mathbf{x}\)，其期望映射为 \(\mathcal{H}(\mathbf{x})\)，则残差块输出可表示为
\begin{equation}
\mathbf{y}=\mathcal{F}(\mathbf{x})+\mathbf{x},
\label{eq:residual_block_ch3}
\end{equation}
其中，\(\mathcal{F}(\mathbf{x})=\mathcal{H}(\mathbf{x})-\mathbf{x}\) 为残差映射。通过式~\eqref{eq:residual_block_ch3} 所示结构，梯度能够沿恒等路径稳定传播，从而缓解深层网络训练中的退化问题。

针对本章任务，网络输入为由式~\eqref{eq:iq_matrix_ch3} 构造的 \(2\times N_s\) 双通道矩阵。前端卷积层首先完成局部波形特征提取与初步下采样，随后通过四组残差阶段进行层级化特征抽取。本文采用的 ResNet-50 主干网络残差阶段配置为 \(\{3,4,6,3\}\)，即四组 Bottleneck 残差块分别包含 3、4、6、3 个单元。随着网络深度增加，浅层卷积主要负责提取局部幅度波动和短时相位变化，中间层逐步建模多径尾迹和跨时延结构，深层网络则进一步抽象出与空间位置相关的高维表示。

经过主干网络特征提取后，本文采用全局平均池化层对空间维度进行压缩，得到共享特征向量 \(\mathbf{f}_{\mathrm{RF}}\in \mathbb{R}^{2048}\)。该特征向量不对应某个单独的传播参数，而是网络对复杂传播环境下 I/Q 信号的幅度、相位和时延结构进行联合编码后的隐式表示。相较于人工提取的低维几何特征，这种高维射频指纹能够更完整地保留位置相关信息，因此更适合复杂环境下的单目标高精度定位任务。

\subsection{损失函数与训练策略}
\label{subsec:3_3_2}

考虑到本章既关注目标所在区域的判别能力，也关注其二维平面坐标的连续估计精度，因此采用分类与回归联合优化的多任务训练策略。设总损失函数为
\begin{equation}
\mathcal{L}_{\mathrm{total}}=\lambda_1\mathcal{L}_{\mathrm{cls}}+\lambda_2\mathcal{L}_{\mathrm{reg}},
\label{eq:total_loss_ch3}
\end{equation}
其中，\(\mathcal{L}_{\mathrm{cls}}\) 为分类损失，\(\mathcal{L}_{\mathrm{reg}}\) 为回归损失，\(\lambda_1\) 与 \(\lambda_2\) 为平衡系数。

分类分支采用 Softmax 输出，并以交叉熵损失作为优化目标：
\begin{equation}
\mathcal{L}_{\mathrm{cls}}=-\frac{1}{M}\sum_{i=1}^{M}\sum_{c=1}^{C}y_{i,c}\log \hat{p}_{i,c},
\label{eq:cls_loss_ch3}
\end{equation}
其中，\(M\) 为批次样本数，\(C\) 为区域类别数，\(y_{i,c}\) 为真实类别标签，\(\hat{p}_{i,c}\) 为预测概率。

回归分支采用均方误差（MSE）损失约束位置预测结果：
\begin{equation}
\mathcal{L}_{\mathrm{reg}}=\frac{1}{M}\sum_{i=1}^{M}\left\|\mathbf{y}^{(\mathrm{reg})}_i-\hat{\mathbf{y}}^{(\mathrm{reg})}_i\right\|_2^2,
\label{eq:reg_loss_ch3}
\end{equation}
其中，\(\mathbf{y}^{(\mathrm{reg})}_i=[x_i,y_i]^T\) 为真实位置标签，\(\hat{\mathbf{y}}^{(\mathrm{reg})}_i\) 为预测坐标。

多任务联合损失的引入具有两方面作用。一方面，分类任务能够增强不同区域样本之间的可分性，使网络优先学习区域层级的空间结构；另一方面，回归任务进一步约束区域内部的位置连续性，从而提高坐标预测精度。通过联合优化，网络可以在“先区分大致区域、再精细回归坐标”的思路下建立更稳定的位置映射关系。

本文采用 Adam 优化器对网络参数进行更新，并结合学习率衰减策略提高收敛稳定性。训练过程中使用验证集监控模型性能，以避免过拟合对测试性能造成影响。本文采用的 ResNet-50 网络结构及训练配置如表~\ref{tab:resnet_params_ch3} 所示。

\begin{table}[htbp]
    \centering
    \bicaption{ResNet-50 射频指纹定位模型结构与训练参数设置}{Model structure and training parameters of ResNet-50 for RF fingerprinting}
    \label{tab:resnet_params_ch3}
    \begin{tabular}{@{}lll@{}}
        \toprule
        \textbf{参数类别} & \textbf{参数名称} & \textbf{取值/设置} \\ \midrule
        输入参数 & 输入形式 & I/Q 双通道矩阵 \\
        & 输入尺寸 & \(2\times N_s\)（\(N_s\) 为时域采样点数） \\
        主干网络 & Backbone & ResNet-50 \\
        & 残差阶段配置 & \(\{3,4,6,3\}\) \\
        & 池化方式 & Global Average Pooling \\
        & 输出特征维度 & 2048 \\
        分支任务 & 分类分支输出数 & \(C\)（由区域划分数确定） \\
        & 回归分支输出维度 & 2（对应平面坐标 \((x,y)\)） \\
        损失函数 & 分类损失函数 & Cross-Entropy \\
        & 回归损失函数 & MSE \\
        & 联合损失形式 & \(\lambda_1 \mathcal{L}_{\mathrm{cls}}+\lambda_2 \mathcal{L}_{\mathrm{reg}}\) \\
        优化策略 & 优化器 & Adam \\
        & Batch Size & 128 \\
        & Epoch & 150 \\
        & 学习率策略 & 衰减策略 \\
        数据集划分 & Train:Val:Test & 8:1:1 \\
        评价指标 & 性能指标 & F1-score, RMSE, Inference Time \\
        扩展支持 & 多天线输入支持 & 支持（可扩展至 MIMO 多通道） \\ \bottomrule
    \end{tabular}
\end{table}

\section{MIMO 空间分集增强机制}
\label{sec:3_4}

在实际非法终端排查场景中，建筑遮挡、传播损耗以及环境噪声会导致接收端处于低信噪比条件。特别是在严重遮挡区域或非视距深谷区域，单天线条件下采集到的 I/Q 信号可能出现明显衰减，进而削弱射频指纹特征的可分性。当输入样本中的有效特征被噪声和深衰落掩盖时，网络虽然仍可输出位置估计结果，但其误差往往会显著增加。因此，仅依赖 SISO 体制难以在极端低信噪比条件下保持稳定性能。

为提升系统在恶劣场景下的鲁棒性，本文在接收端引入多天线空间分集机制。设接收端配置 \(K\) 个天线单元，则同一时刻可获得多路空间观测，构成多通道接收矩阵
\begin{equation}
\mathbf{R}_{\mathrm{MIMO}}[n]=
\begin{bmatrix}
r_1[n] & r_2[n] & \cdots & r_K[n]
\end{bmatrix}^{T},
\label{eq:mimo_signal_ch3}
\end{equation}
其中，\(r_k[n]\) 表示第 \(k\) 个天线通道对应的复采样序列。

多天线配置的作用并不在于显式实现传统波束形成或最大比合并算法，而在于为网络提供更丰富的空间观测信息。由于不同天线单元所处空间位置不同，各通道对多径结构和局部深衰落的响应并不完全一致。当某一通道受到严重衰落影响时，其他通道仍可能保留较强的有效观测，从而在整体上提高输入样本的稳定性。基于这一思路，本文将各天线通道的 I/Q 样本沿通道维进行拼接，使原始双通道输入扩展为 \(2K\) 通道输入张量，再送入深度网络进行联合特征学习。相较于单通道输入，该方式能够更好地保留低信噪比条件下的空间分集信息，从而提高模型对目标位置的判别能力。

本节的核心在于说明：MIMO 机制并不替代深度射频指纹学习，而是为其提供更稳定的输入基础。后续实验将通过不同阵列规模与不同 SNR 条件下的性能对比，验证空间分集对复杂环境单目标定位精度和鲁棒性的提升作用。

\section{实验设置与评价指标}
\label{sec:3_5}

为验证所提方法在复杂城市环境中的单目标定位性能，本文构建了基于三维城市场景建模与射线追踪仿真的实验平台。实验区域采用典型建筑密集场景，数据集覆盖目标区域内 27,546 个网格位置，并在不同信噪比和不同接收阵列配置条件下生成相应的 I/Q 指纹样本。信号载频设置在低轨卫星相关 Ku 波段范围内，噪声模型采用复加性高斯白噪声，通过改变 SNR 构造从中高信噪比到极低信噪比的不同测试条件。

在接收阵列配置方面，本文分别考虑 SISO、\(2\times2\)、\(4\times4\) 和 \(8\times8\) 四种体制，用于分析空间分集对定位性能的影响。所有模型在相同的数据划分和统一训练策略下进行对比，以保证实验结果的可比性。训练集、验证集和测试集按照 8:1:1 比例划分，训练批大小设置为 128，最大训练轮次设置为 150。

在评价指标方面，本文主要采用均方根误差（RMSE）衡量连续坐标预测精度，其定义为
\begin{equation}
\mathrm{RMSE}=\sqrt{\frac{1}{N_{\mathrm{test}}}\sum_{i=1}^{N_{\mathrm{test}}}\left[(x_i-\hat{x}_i)^2+(y_i-\hat{y}_i)^2\right]},
\label{eq:rmse_ch3}
\end{equation}
其中，\(N_{\mathrm{test}}\) 为测试样本数。对于区域分类任务，采用 F1-score 作为综合评价指标；对于工程部署能力，则进一步统计模型单次前向推理时间，以评估其在边缘平台上的应用潜力。

\section{实验验证与性能分析}
\label{sec:3_6}

\subsection{定位精度对比分析}
\label{subsec:3_6_1}

为比较不同方法在复杂城市环境下的单目标定位性能，本文将几何特征方法、浅层卷积网络和深层残差网络的测试结果统一列于表~\ref{tab:perf_comparison_ch3}。

\begin{table}[htbp]
    \centering
    \bicaption{不同单目标定位方法性能对比}{Performance comparison of different single-target localization methods}
    \label{tab:perf_comparison_ch3}
    \begin{tabular}{@{}llcccc@{}}
        \toprule
        \textbf{方法} & \textbf{输入特征} & \textbf{模型类型} & \textbf{RMSE / m} & \textbf{F1-score} & \textbf{推理时间 / ms} \\ \midrule
        Geometric-based & RSSI、AoA 等低维几何特征 & 传统回归模型 & 45.2 & -- & -- \\
        4-Layer CNN & I/Q 样本 & 浅层卷积网络 & 待补 & 0.73 & -- \\
        ResNet-50 & I/Q 样本 & 深层残差网络 & \textbf{8.3} & \textbf{0.85} & \textbf{85.0} \\
        ResNet-101 & I/Q 样本 & 更深层残差网络 & 待补 & 待补 & 142.0 \\ \bottomrule
    \end{tabular}
\end{table}

由表~\ref{tab:perf_comparison_ch3} 可见，基于低维几何特征的传统方法 RMSE 约为 \(45.2\,\mathrm{m}\)，说明在复杂多径环境中仅依赖少量观测量难以获得稳定的高精度定位结果。浅层 4 层卷积网络在区域分类任务中可取得 0.73 的 F1-score，而本文采用的 ResNet-50 方法可将定位 RMSE 降低至 \(8.3\,\mathrm{m}\)，对应区域分类 F1-score 提升至 0.85。该结果表明，直接从原始 I/Q 复基带样本中提取深层射频指纹，相较于依赖低维人工特征或浅层网络的方法，更有利于恢复与位置相关的传播结构信息。

从原因上看，几何特征方法在输入阶段就丢失了大量与多径传播相关的细粒度信息，因此面对复杂环境时难以准确描述位置差异。浅层 CNN 虽然具备一定局部特征提取能力，但对长跨度时延结构和多路径相位耦合的建模能力仍然有限。相比之下，ResNet-50 借助更强的深层表达能力和残差连接机制，能够在较大感受野内联合建模复杂传播特征，因此获得更低的定位误差。

\subsection{计算开销与鲁棒性折衷}
\label{subsec:3_6_2}

为进一步分析网络深度对性能与计算代价的影响，本文对 ResNet-50 与更深层的 ResNet-101 进行了比较。实验结果表明，随着网络深度增加，模型的表达能力可能有所增强，但计算代价也显著提高。对于本章任务而言，ResNet-50 在精度和推理时延之间取得了较好的平衡，而更深的 ResNet-101 虽然可能带来有限性能提升，但单次前向推理时间增加至 \(142.0\,\mathrm{ms}\)，对低成本无人机平台而言部署压力更大。

相比之下，ResNet-50 的单次前向推理时间约为 \(85.0\,\mathrm{ms}\)。在保证较高定位精度的同时，该时延水平更适合作为边缘平台上的主干网络。由此可见，单纯增加网络深度并不一定能够显著改善复杂传播环境下的定位性能，而合理控制网络结构复杂度对于工程部署同样重要。

\subsection{极端信噪比条件下的鲁棒性评估}
\label{subsec:3_6_3}

在复杂城市环境中，建筑遮挡和传播损耗可能导致接收端处于极低信噪比条件。为验证所提方法在恶劣条件下的适用性，本文进一步构造了从中高信噪比逐步下降至 \(-15\,\mathrm{dB}\) 的测试场景，并分析不同接收体制下模型性能随 SNR 的变化趋势。

实验表明，随着 SNR 持续下降，各类方法的定位误差均有所增加，但不同阵列配置下的误差增长速度存在明显差异。单天线条件下，接收信号更容易受到局部深衰落和噪声污染影响，因此模型性能退化较快；而在引入空间分集后，多通道输入能够在一定程度上保留更稳定的传播特征，使得模型在低信噪比区域仍保持一定辨识能力。这说明，在复杂环境中，网络性能不仅取决于结构设计，也与输入观测的稳定性密切相关。

\subsection{不同 MIMO 配置下的性能分析}
\label{subsec:3_6_4}

为进一步分析空间分集机制对低信噪比场景定位性能的影响，本文将不同接收阵列配置下的实验结果列于表~\ref{tab:mimo_comparison_ch3}。

\begin{table}[htbp]
    \centering
    \bicaption{不同 MIMO 配置下低信噪比场景定位性能对比}{Localization performance with different MIMO configurations in low SNR scenarios}
    \label{tab:mimo_comparison_ch3}
    \begin{tabular}{@{}lcccl@{}}
        \toprule
        \textbf{接收配置} & \textbf{输入通道数} & \textbf{典型 SNR 条件} & \textbf{RMSE / m} & \textbf{性能描述} \\ \midrule
        SISO & 2 & \(-15\,\mathrm{dB}\) & 待补 & 误差显著增大 \\
        \(2\times 2\) MIMO & 4 & \(-15\,\mathrm{dB}\) & 待补 & 较 SISO 有明显改善 \\
        \(4\times 4\) MIMO & 8 & \(-15\,\mathrm{dB}\) & 待补 & 分集进一步降低误差 \\
        \(8\times 8\) MIMO & 16 & \(-15\,\mathrm{dB}\) & \(\mathbf{<20}\) & 保持较好鲁棒性 \\ \bottomrule
    \end{tabular}
\end{table}

表~\ref{tab:mimo_comparison_ch3} 表明，在极低信噪比条件下，不同接收阵列配置的定位性能差异较为明显。SISO 体制下，系统仅能依赖单通道观测，输入特征受噪声和深衰落影响较大，因而定位误差显著增大。随着阵列规模增加，多通道空间分集能够在更大程度上保留目标相关特征，从而降低噪声和局部深衰落对定位结果的影响。特别是在 \(-15\,\mathrm{dB}\) 条件下，\(8\times8\) MIMO 仍能将平均定位误差控制在 \(20\,\mathrm{m}\) 以内，说明多天线空间分集对复杂环境下射频指纹可分性的提升具有积极作用。

需要指出的是，阵列规模增加虽然有利于提升鲁棒性，但也会带来更高的数据维度和计算负担。因此，在实际部署中仍需结合平台载荷和算力条件对阵列规模进行权衡。总体而言，本节结果验证了 MIMO 空间分集机制对低信噪比复杂场景定位性能的增强作用。

\section{本章小结}
\label{sec:3_7}

本章面向复杂城市环境中的单目标定位问题，提出了一种基于深度射频指纹的高精度定位方法。首先，结合 OSM 城市建筑数据和高保真射线追踪技术，构建了复杂传播场景下的位置相关样本生成框架；其次，基于低轨卫星相关下行链路参考信号，生成包含多径传播特征的 I/Q 双路复基带样本，并将其作为深度射频指纹学习的原始输入；随后，采用 ResNet-50 作为主干网络，通过区域分类与位置回归联合优化实现单目标定位；最后，通过引入多天线空间分集机制，增强了系统在低信噪比条件下的鲁棒性。

实验结果表明，相较于基于低维几何特征的传统方法，所提方法能够更有效地利用复杂传播环境中的位置相关信息，从而在复杂城市环境下获得更低的定位误差。与此同时，ResNet-50 在定位精度与推理效率之间取得了较好平衡，而多天线输入进一步提升了系统在恶劣传播条件下的特征稳定性。上述研究为后续多目标定位与多源信息融合方法提供了较高质量的单目标特征基础。